\documentclass[10pt,a4paper]{amsart}
\usepackage{amsmath,amssymb,amsthm,graphicx}
\usepackage{verbatim}
\newenvironment{reference}{\comment}{\endcomment}
%\newenvironment{reference}{}{}
\newenvironment{slogan}{\comment}{\endcomment}
\newenvironment{history}{\comment}{\endcomment}

% For commutative diagrams we use Xy-pic
\usepackage[all]{xy}

% We use 2cell for 2-commutative diagrams.
\xyoption{2cell}
% Theorem environments.
%
\theoremstyle{plain}
\newtheorem{theorem}[subsection]{Theorem}
\newtheorem{proposition}[subsection]{Proposition}
\newtheorem{lemma}[subsection]{Lemma}

\theoremstyle{definition}
\newtheorem{definition}[subsection]{Definition}
\newtheorem{example}[subsection]{Example}
\newtheorem{exercise}[subsection]{Exercise}
\newtheorem{situation}[subsection]{Situation}

\theoremstyle{remark}
\newtheorem{remark}[subsection]{Remark}
\newtheorem{remarks}[subsection]{Remarks}

\numberwithin{equation}{subsection}

% Macros
%
\def\lim{\mathop{\mathrm{lim}}\nolimits}
\def\colim{\mathop{\mathrm{colim}}\nolimits}
\def\Spec{\mathop{\mathrm{Spec}}}
\def\Hom{\mathop{\mathrm{Hom}}\nolimits}
\def\Ext{\mathop{\mathrm{Ext}}\nolimits}
\def\SheafHom{\mathop{\mathcal{H}\!\mathit{om}}\nolimits}
\def\SheafExt{\mathop{\mathcal{E}\!\mathit{xt}}\nolimits}
\def\Sch{\mathit{Sch}}
\def\Mor{\mathop{\mathrm{Mor}}\nolimits}
\def\Ob{\mathop{\mathrm{Ob}}\nolimits}
\def\Sh{\mathop{\mathit{Sh}}\nolimits}
\def\NL{\mathop{N\!L}\nolimits}
\def\CH{\mathop{\mathrm{CH}}\nolimits}
\def\proetale{{pro\text{-}\acute{e}tale}}
\def\etale{{\acute{e}tale}}
\def\QCoh{\mathit{QCoh}}
\def\Ker{\mathop{\mathrm{Ker}}}
\def\Im{\mathop{\mathrm{Im}}}
\def\Coker{\mathop{\mathrm{Coker}}}
\def\Coim{\mathop{\mathrm{Coim}}}

% Boxtimes
%
\DeclareMathSymbol{\boxtimes}{\mathbin}{AMSa}{"02}

%
% Macros for moduli stacks/spaces
%
\def\QCohstack{\mathcal{QC}\!\mathit{oh}}
\def\Cohstack{\mathcal{C}\!\mathit{oh}}
\def\Spacesstack{\mathcal{S}\!\mathit{paces}}
\def\Quotfunctor{\mathrm{Quot}}
\def\Hilbfunctor{\mathrm{Hilb}}
\def\Curvesstack{\mathcal{C}\!\mathit{urves}}
\def\Polarizedstack{\mathcal{P}\!\mathit{olarized}}
\def\Complexesstack{\mathcal{C}\!\mathit{omplexes}}
% \Pic is the operator that assigns to X its picard group, usage \Pic(X)
% \Picardstack_{X/B} denotes the Picard stack of X over B
% \Picardfunctor_{X/B} denotes the Picard functor of X over B
\def\Pic{\mathop{\mathrm{Pic}}\nolimits}
\def\Picardstack{\mathcal{P}\!\mathit{ic}}
\def\Picardfunctor{\mathrm{Pic}}
\def\Deformationcategory{\mathcal{D}\!\mathit{ef}}
\usepackage{cleveref}
\crefname{originalsection}{original section}{original sections}
\Crefname{originalsection}{Original section}{Original sections}
\makeatletter
\newlabel{section-cycle-of-closed}{{0AZA}{source}{}{sourcepointer-1}{}}
\newlabel{section-cycle-of-closed@cref}{{[originalsection][0][]0AZA}{[1][0][]source}{}{}{}}
\makeatother
\crefname{originalreference}{original reference}{original references}
\Crefname{originalreference}{Original reference}{Original references}
\makeatletter
\newlabel{varieties-lemma-dimension-locally-algebraic}{{0A21}{source}{}{sourcepointer-2}{}}
\newlabel{varieties-lemma-dimension-locally-algebraic@cref}{{[originalreference][0][]0A21}{[1][0][]source}{}{}{}}
\makeatother
\crefname{originalreference}{original reference}{original references}
\Crefname{originalreference}{Original reference}{Original references}
\makeatletter
\newlabel{algebra-definition-regular-local}{{00KU}{source}{}{sourcepointer-3}{}}
\newlabel{algebra-definition-regular-local@cref}{{[originalreference][0][]00KU}{[1][0][]source}{}{}{}}
\makeatother
\crefname{originalreference}{original reference}{original references}
\Crefname{originalreference}{Original reference}{Original references}
\makeatletter
\newlabel{algebra-lemma-regular-ring-CM}{{00NQ}{source}{}{sourcepointer-4}{}}
\newlabel{algebra-lemma-regular-ring-CM@cref}{{[originalreference][0][]00NQ}{[1][0][]source}{}{}{}}
\makeatother
\crefname{originalreference}{original reference}{original references}
\Crefname{originalreference}{Original reference}{Original references}
\makeatletter
\newlabel{algebra-lemma-regular-domain}{{00NP}{source}{}{sourcepointer-5}{}}
\newlabel{algebra-lemma-regular-domain@cref}{{[originalreference][0][]00NP}{[1][0][]source}{}{}{}}
\makeatother
\crefname{originalreference}{original reference}{original references}
\Crefname{originalreference}{Original reference}{Original references}
\makeatletter
\newlabel{algebra-lemma-one-equation}{{00KW}{source}{}{sourcepointer-6}{}}
\newlabel{algebra-lemma-one-equation@cref}{{[originalreference][0][]00KW}{[1][0][]source}{}{}{}}
\makeatother
\crefname{originalreference}{original reference}{original references}
\Crefname{originalreference}{Original reference}{Original references}
\makeatletter
\newlabel{more-algebra-lemma-regular-koszul-regular}{{062F}{source}{}{sourcepointer-7}{}}
\newlabel{more-algebra-lemma-regular-koszul-regular@cref}{{[originalreference][0][]062F}{[1][0][]source}{}{}{}}
\makeatother
\title{Multiplicity one forces transversal intersection: e = 1 and smoothness in general points}
\author{The Stacks Project Authors}
\date{stacks-project:intersection.tex:local-snapshot-20261012; GNU FDL 1.2 or later}
\begin{document}
\maketitle
\noindent\textit{Editorial note.} All delivered mathematical text is the author's own, copied verbatim from the pinned Stacks Project chapter file intersection.tex; no statement or proof step was rewritten, reordered or removed.\par\medskip
\noindent\textit{Editorial note.} Span recomputation by label pairing (LANE\_BRIEF 10.2): in this chapter the author's proof environment follows the statement environment as a standalone block; the spec covers the complete environment pair (statement span, proof span) as recorded in the pool candidate.\par\medskip
\noindent\textit{Editorial note.} Occupation scan: combined claim\_key + statement\_anchor + generic-line scan over all registration files (active and withdrawn) found 0 hits for this chapter (is-occupied.json, 2026-10-14); consistent with the duty manager's six-check dispatch.\par\medskip
\ifdefined\hypertarget\hypertarget{sourcepointer-1}{}\fi
\noindent\textit{Source reference.} Original section 0AZA: sources/stacks-chow/tags.txt, source lines 13818--13818. Cross-chapter standard prerequisite cited by the author inside the counted span; it is used as a pointer and its own proof is not part of this item.\par\medskip
\ifdefined\hypertarget\hypertarget{sourcepointer-2}{}\fi
\noindent\textit{Source reference.} Original reference 0A21: sources/stacks-chow/tags.txt, source lines 12689--12689. Cross-chapter standard prerequisite cited by the author inside the counted span; it is used as a pointer and its own proof is not part of this item.\par\medskip
\ifdefined\hypertarget\hypertarget{sourcepointer-3}{}\fi
\noindent\textit{Source reference.} Original reference 00KU: sources/stacks-chow/tags.txt, source lines 764--764. Cross-chapter standard prerequisite cited by the author inside the counted span; it is used as a pointer and its own proof is not part of this item.\par\medskip
\ifdefined\hypertarget\hypertarget{sourcepointer-4}{}\fi
\noindent\textit{Source reference.} Original reference 00NQ: sources/stacks-chow/tags.txt, source lines 868--868. Cross-chapter standard prerequisite cited by the author inside the counted span; it is used as a pointer and its own proof is not part of this item.\par\medskip
\ifdefined\hypertarget\hypertarget{sourcepointer-5}{}\fi
\noindent\textit{Source reference.} Original reference 00NP: sources/stacks-chow/tags.txt, source lines 867--867. Cross-chapter standard prerequisite cited by the author inside the counted span; it is used as a pointer and its own proof is not part of this item.\par\medskip
\ifdefined\hypertarget\hypertarget{sourcepointer-6}{}\fi
\noindent\textit{Source reference.} Original reference 00KW: sources/stacks-chow/tags.txt, source lines 766--766. Cross-chapter standard prerequisite cited by the author inside the counted span; it is used as a pointer and its own proof is not part of this item.\par\medskip
\ifdefined\hypertarget\hypertarget{sourcepointer-7}{}\fi
\noindent\textit{Source reference.} Original reference 062F: sources/stacks-chow/tags.txt, source lines 7794--7794. Cross-chapter standard prerequisite cited by the author inside the counted span; it is used as a pointer and its own proof is not part of this item.\par\medskip
\begin{lemma}
\label{lemma-transversal}
Let $X$ be a nonsingular variety. Let $V, W \subset X$ be
closed subvarieties which intersect properly. Let $Z$ be an irreducible
component of $V \cap W$ and assume that the multiplicity
(in the sense of Section \ref{section-cycle-of-closed}) of $Z$
in the closed subscheme $V \cap W$ is $1$.
Then $e(X, V \cdot W, Z) = 1$ and $V$ and $W$ are smooth
in a general point of $Z$.
\end{lemma}

\begin{proof}
Let $(A, \mathfrak m, \kappa) =
(\mathcal{O}_{X, \xi}, \mathfrak m_\xi, \kappa(\xi))$ where $\xi \in Z$
is the generic point. Then $\dim(A) = \dim(X) - \dim(Z)$, see
Varieties, Lemma \ref{varieties-lemma-dimension-locally-algebraic}.
Let $I, J \subset A$ cut out the trace of $V$ and $W$
in $\Spec(A)$. Set $\overline{I} = I + \mathfrak m^2/\mathfrak m^2$.
Then $\dim_\kappa \overline{I} \leq \dim(X) - \dim(V)$ with equality
if and only if $A/I$ is regular (this follows from the lemma cited
above and the definition of regular rings, see
Algebra, Definition \ref{algebra-definition-regular-local}
and the discussion preceding it). Similarly for $\overline{J}$.
If the multiplicity is $1$, then
$\text{length}_A(A/I + J) = 1$, hence $I + J = \mathfrak m$, hence
$\overline{I} + \overline{J} = \mathfrak m/\mathfrak m^2$.
Then we get equality everywhere (because the intersection is
proper). Hence we find $f_1, \ldots, f_a \in I$ and $g_1, \ldots g_b \in J$
such that $\overline{f}_1, \ldots, \overline{g}_b$ is a basis
for $\mathfrak m/\mathfrak m^2$. Then $f_1, \ldots, g_b$ is a
regular system of parameters and a regular sequence
(Algebra, Lemma \ref{algebra-lemma-regular-ring-CM}).
The same lemma shows $A/(f_1, \ldots, f_a)$ is a regular local ring
of dimension $\dim(X) - \dim(V)$, hence $A/(f_1, \ldots, f_a) \to A/I$
is an isomorphism (if the kernel is nonzero, then the dimension
of $A/I$ is strictly less, see
Algebra, Lemmas \ref{algebra-lemma-regular-domain} and
\ref{algebra-lemma-one-equation}).
We conclude $I = (f_1, \ldots, f_a)$ and $J = (g_1, \ldots, g_b)$
by symmetry. Thus the Koszul complex $K_\bullet(A, f_1, \ldots, f_a)$
on $f_1, \ldots, f_a$ is a resolution of $A/I$, see
More on Algebra, Lemma \ref{more-algebra-lemma-regular-koszul-regular}.
Hence
\begin{align*}
\text{Tor}_p^A(A/I, A/J)
& =
H_p(K_\bullet(A, f_1, \ldots, f_a) \otimes_A A/J) \\
& =
H_p(K_\bullet(A/J, f_1 \bmod J, \ldots, f_a \bmod J))
\end{align*}
Since we've seen above that $f_1 \bmod J, \ldots, f_a \bmod J$ is
a regular system of parameters in the regular local ring $A/J$
we conclude that there is only one cohomology group, namely
$H_0 = A/(I + J) = \kappa$. This finishes the proof.
\end{proof}

\end{document}
